% =============================================================================
% Presentación: Artificial Neural Networks for The Perceptron, Madaline,
%   and Backpropagation Family (Bernard Widrow & Michael A. Lehr)
% Presentado por: Zhiyong Cheng — 2014-09-08
% =============================================================================
\documentclass[aspectratio=169,utf8,spanish]{beamer}

% ── Paquetes ────────────────────────────────────────────────────────────────
\usepackage[utf8]{inputenc}
\usepackage[T1]{fontenc}
\usepackage[spanish,es-noshorthands]{babel}
\usepackage{amsmath,amssymb,amsthm}
\usepackage{graphicx}
\usepackage{booktabs}
\usepackage{hyperref}
\usepackage{xcolor}

% ── Tema ────────────────────────────────────────────────────────────────────
\usetheme{Madrid}
\usecolortheme{default}

% ── Metadatos ───────────────────────────────────────────────────────────────
\title[ANN: Perceptron, Madaline, Backpropagation]{
  Artificial Neural Networks for The Perceptron,\\
  Madaline, and Backpropagation Family
}
\subtitle{Bernard Widrow \& Michael A. Lehr}
\author{Presentado por Zhiyong Cheng}
\institute{}
\date{2014-09-08}

% ── Atajos ──────────────────────────────────────────────────────────────────
\newcommand{\R}{\mathbb{R}}
\newcommand{\E}{\mathbb{E}}
\DeclareMathOperator{\sgn}{sgn}
\DeclareMathOperator{\sgm}{sgm}

\begin{document}

% ── Portada ─────────────────────────────────────────────────────────────────
\begin{frame}
  \titlepage
\end{frame}

% ═══════════════════════════════════════════════════════════════════════════
\section{Contenido}
% ═══════════════════════════════════════════════════════════════════════════

\begin{frame}{Contenido (Outline)}
  \begin{itemize}
    \item Historia del desarrollo de algoritmos.
    \item Conceptos fundamentales.
    \item Algoritmos:
    \begin{itemize}
      \item Reglas de corrección de error:
      \begin{itemize}
        \item Elemento único: lineal y no lineal.
        \item Redes en capas: no lineal.
      \end{itemize}
      \item Descenso de gradiente:
      \begin{itemize}
        \item Elemento único: lineal y no lineal.
        \item Redes en capas: no lineal.
      \end{itemize}
    \end{itemize}
  \end{itemize}
\end{frame}

% ═══════════════════════════════════════════════════════════════════════════
\section{Evolución de Algoritmos}
% ═══════════════════════════════════════════════════════════════════════════

\begin{frame}{Evolución de Algoritmos}
  \begin{itemize}
    \item \textbf{1960:} Dos de las reglas más tempranas e importantes para entrenar elementos adaptativos:
    \begin{itemize}
      \item Algoritmo LMS (Widrow y su estudiante).
      \item Regla del Perceptrón (Rosenblatt).
    \end{itemize}
    \item Facilitaron el rápido desarrollo de redes neuronales en los primeros años:
    \begin{itemize}
      \item Madaline Rule I (MRI): primera regla de aprendizaje popular para redes neuronales con múltiples elementos adaptativos.
      \item Aplicaciones: reconocimiento de voz y patrones, predicción del clima, etc.
    \end{itemize}
    \item \textbf{Avance decisivo:} algoritmo de entrenamiento por retro-propagación (\textit{backpropagation}):
    \begin{itemize}
      \item Desarrollado por Werbos en 1971, publicado en 1974.
      \item Redescubierto por Parker en 1982, reportado en 1984.
      \item Rumelhart, Hinton y Williams también lo redescubrieron y lo hicieron ampliamente conocido.
    \end{itemize}
  \end{itemize}
\end{frame}

\begin{frame}{Evolución de Algoritmos (cont.)}
  \begin{itemize}
    \item \textbf{MRI:} cuantizadores duros (\textit{hard-limiting}, $y = \sgn(x)$); fácil de calcular.
    \item \textbf{Backpropagation:} función sigmoide.
    \item \textbf{MRI $\to$ MRII} (Widrow y Winter, 1987).
    \item \textbf{MRII $\to$ MRIII} (David Andes, U.S. Naval Weapons Center, China Lake, CA, 1988):
    \begin{itemize}
      \item Reemplaza cuantizadores duros por función sigmoide.
      \item Widrow y su estudiante descubren primero que MRIII es equivalente a backpropagation.
    \end{itemize}
  \end{itemize}
\end{frame}

% ═══════════════════════════════════════════════════════════════════════════
\section{Conceptos Básicos}
% ═══════════════════════════════════════════════════════════════════════════

\begin{frame}{Combinador Lineal Adaptativo}
  \begin{columns}
    \begin{column}{0.5\textwidth}
      \begin{itemize}
        \item Entradas: $x_1, x_2, x_3$ (más $x_0$, sesgo).
        \item Salida: $y = \mathbf{X}_k^T \mathbf{W}_k$.
        \item Capa de entrada $\to$ Capa de salida.
      \end{itemize}
    \end{column}
    \begin{column}{0.5\textwidth}
      \begin{center}
        \begin{tabular}{c}
          $x_0$ \\
          $x_1$ \\
          $x_2$ \\
          $x_3$
        \end{tabular}
        $\to$
        $\displaystyle\sum$
        $\to y$
      \end{center}
    \end{column}
  \end{columns}
\end{frame}

\begin{frame}{Elemento Lineal Adaptativo (Adaline) — Clasificador Lineal}
  \begin{columns}
    \begin{column}{0.5\textwidth}
      \begin{itemize}
        \item Combinador lineal + cuantizador signum.
        \item $y = \sgn(\mathbf{X}_k^T \mathbf{W}_k)$.
        \item Separa el espacio mediante un hiperplano.
      \end{itemize}
    \end{column}
    \begin{column}{0.5\textwidth}
      \begin{center}
        \begin{tabular}{c}
          $x_0$ \\
          $x_1$ \\
          $x_2$ \\
          $x_3$
        \end{tabular}
        $\to$
        $\displaystyle\sum$
        $\to \sgn(\cdot) \to y$
      \end{center}
    \end{column}
  \end{columns}

  \begin{center}
    \textit{Figuras tomadas del curso online de Andrew Ng}
  \end{center}
\end{frame}

\begin{frame}{Clasificadores No Lineales}
  \begin{itemize}
    \item \textbf{Funciones discriminantes polinomiales.}
    \item \textbf{Red neuronal feedforward multi-elemento.}
  \end{itemize}
  \begin{center}
    \textit{Figuras tomadas del curso online de Andrew Ng}
  \end{center}
\end{frame}

\begin{frame}{Función Discriminante Polinomial}
  \begin{itemize}
    \item Vector de características inicial: $[x_1, x_2]$.
    \item Después del preprocesador: $[x_1^2,\; x_1,\; x_1x_2,\; x_2,\; x_2^2]$.
    \item Nuevo vector de características: $[x_1', x_2', x_3', x_4', x_5']$.
  \end{itemize}

  \begin{block}{Ventajas}
    \begin{itemize}
      \item Puede proporcionar una variedad de funciones no lineales con un solo elemento Adaline.
      \item Converge más rápido que las redes neuronales en capas.
      \item Solución global única.
    \end{itemize}
  \end{block}

  \begin{block}{Desventaja}
    Las redes neuronales en capas pueden obtener mejor generalización.
  \end{block}
\end{frame}

\begin{frame}{Redes Neuronales en Capas: Madaline I}
  \begin{itemize}
    \item Múltiples elementos Adaline en la primera capa.
    \item Dispositivos lógicos fijos en la segunda capa (OR, AND, voto mayoritario, etc.).
  \end{itemize}
  \begin{center}
    $\mathbf{X} \to$ [Adaline$_1$, $\ldots$, Adaline$_n$] $\to$ Lógica Fija $\to \pm 1$

    \small{\textit{Primera capa} \hspace{3cm} \textit{Segunda capa}}
  \end{center}
\end{frame}

\begin{frame}{Red Feedforward}
  \begin{itemize}
    \item Todas las capas son adaptativas.
    \item Múltiples capas ocultas entre entrada y salida.
  \end{itemize}
  \begin{center}
    Capa 1 (entrada) $\to$ Capa 2 (oculta) $\to$ Capa 3 (oculta) $\to$ Capa 4 (salida)
  \end{center}
\end{frame}

% ═══════════════════════════════════════════════════════════════════════════
\section{El Principio de Mínima Perturbación}
% ═══════════════════════════════════════════════════════════════════════════

\begin{frame}{El Principio de Mínima Perturbación}
  \begin{block}{Principio de Mínima Perturbación}
    Adaptar para reducir el error de salida para el patrón de entrenamiento actual, con mínima perturbación a las respuestas ya aprendidas.
  \end{block}

  \begin{itemize}
    \item Condujo al descubrimiento del algoritmo LMS y las reglas Madaline.
    \item \textbf{Adaline único:} $\alpha$-LMS, $\mu$-LMS, reglas de May, regla del Perceptrón.
    \item \textbf{Madaline simple:} MRI.
    \item \textbf{Madaline multi-capa:} MRII, MRIII y algoritmo de backpropagation.
  \end{itemize}
\end{frame}

% ═══════════════════════════════════════════════════════════════════════════
\section{Algoritmos Adaptativos}
% ═══════════════════════════════════════════════════════════════════════════

\begin{frame}{Algoritmos Adaptativos: Dos Familias}
  \begin{block}{Reglas de Corrección de Error}
    Alteran los pesos para corregir una cierta proporción del error en la respuesta de salida al patrón de entrada presente.
    \begin{itemize}
      \item $\alpha$-LMS: $\mathbf{W}_{k+1} = \mathbf{W}_k + \alpha \dfrac{\varepsilon_k \mathbf{X}_k}{\|\mathbf{X}_k\|^2}$
      \item Corrección de error: $\Delta\varepsilon_k = -\alpha\varepsilon_k$
    \end{itemize}
  \end{block}

  \begin{block}{Reglas de Gradiente (Descenso más Empinado)}
    Alteran los pesos durante cada presentación de patrón por descenso de gradiente con el objetivo de reducir el error cuadrático medio, promediado sobre todos los patrones.
    \begin{itemize}
      \item $\mu$-LMS: $\mathbf{W}_{k+1} = \mathbf{W}_k + 2\mu\varepsilon_k \mathbf{X}_k$
    \end{itemize}
  \end{block}
\end{frame}

% ═══════════════════════════════════════════════════════════════════════════
\section{Reglas de Corrección de Error: Adaline}
% ═══════════════════════════════════════════════════════════════════════════

\begin{frame}{Regla Lineal: $\alpha$-LMS}
  \begin{itemize}
    \item \textbf{Actualización de pesos:}
    $\mathbf{W}_{k+1} = \mathbf{W}_k + \alpha \dfrac{\varepsilon_k \mathbf{X}_k}{\|\mathbf{X}_k\|^2}$
    \item \textbf{Error:} $\varepsilon_k = d_k - \mathbf{W}_k^T \mathbf{X}_k$
    \item \textbf{Cambio de error por el cambio de peso:}
    \begin{align*}
      \Delta\varepsilon_k &= \Delta(d_k - \mathbf{W}_k^T \mathbf{X}_k) = -\mathbf{X}_k^T \Delta\mathbf{W}_k \\
      \Delta\mathbf{W}_k &= \mathbf{W}_{k+1} - \mathbf{W}_k = \alpha \frac{\varepsilon_k \mathbf{X}_k}{\|\mathbf{X}_k\|^2} \\
      \Delta\varepsilon_k &= -\alpha \frac{\varepsilon_k \mathbf{X}_k^T \mathbf{X}_k}{\|\mathbf{X}_k\|^2} = -\alpha\varepsilon_k
    \end{align*}
    \item \textbf{Tasa de aprendizaje:} $0.1 < \alpha < 1$.
  \end{itemize}
\end{frame}

\begin{frame}{Regla No Lineal: Perceptrón}
  \begin{itemize}
    \item \textbf{Error del cuantizador:} $\tilde{\varepsilon}_k = d_k - y_k$ (donde $y_k = \sgn(s_k)$).
    \item \textbf{Actualización de pesos:}
    $\mathbf{W}_{k+1} = \mathbf{W}_k + \alpha \dfrac{\tilde{\varepsilon}_k}{2} \mathbf{X}_k$
    \item Para $\alpha = 1$:
    \begin{itemize}
      \item Si $\tilde{\varepsilon}_k > 0$ (i.e., $d_k = 1$, $y_k = -1$): $\mathbf{W}_{k+1} = \mathbf{W}_k + \mathbf{X}_k$.
      \item Si $\tilde{\varepsilon}_k < 0$ (i.e., $d_k = -1$, $y_k = 1$): $\mathbf{W}_{k+1} = \mathbf{W}_k - \mathbf{X}_k$.
    \end{itemize}
  \end{itemize}
\end{frame}

\begin{frame}{$\alpha$-LMS vs Regla del Perceptrón}
  \begin{itemize}
    \item \textbf{Valor de $\alpha$:}
    \begin{itemize}
      \item $\alpha$-LMS: controla estabilidad y velocidad de convergencia.
      \item Perceptrón: no afecta la estabilidad; solo afecta el tiempo de convergencia si el vector de pesos inicial es no nulo.
    \end{itemize}
    \item \textbf{Tipo de respuesta:}
    \begin{itemize}
      \item $\alpha$-LMS: respuestas binarias y continuas.
      \item Perceptrón: solo respuestas binarias deseadas.
    \end{itemize}
    \item \textbf{Patrones linealmente separables:}
    \begin{itemize}
      \item Perceptrón: separa cualquier conjunto linealmente separable.
      \item $\alpha$-LMS: puede fallar en separar conjuntos linealmente separables.
    \end{itemize}
    \item \textbf{Patrones no linealmente separables:}
    \begin{itemize}
      \item Perceptrón: no converge; a menudo no produce una solución de bajo error.
      \item $\alpha$-LMS: no conduce a soluciones de peso no razonables.
    \end{itemize}
  \end{itemize}
\end{frame}

\begin{frame}{Regla No Lineal: Algoritmos de May}
  \begin{itemize}
    \item Separan cualquier conjunto linealmente separable.
    \item Para conjuntos no linealmente separables, la regla de May funciona mucho mejor que el Perceptrón porque una zona muerta suficientemente grande tiende a hacer que el vector de pesos se adapte lejos de cero cuando existe una solución razonablemente buena.
  \end{itemize}

  \begin{block}{Variantes}
    \begin{itemize}
      \item \textbf{Regla de Adaptación por Incremento.}
      \item \textbf{Regla de Adaptación por Relajación Modificada.}
    \end{itemize}
  \end{block}
\end{frame}

% ═══════════════════════════════════════════════════════════════════════════
\section{Reglas de Corrección de Error: Madaline}
% ═══════════════════════════════════════════════════════════════════════════

\begin{frame}{Madaline Rule I (MRI) — Arquitectura}
  \begin{itemize}
    \item \textbf{Primera capa:} elementos Adaline con cuantizador duro.
    \item \textbf{Segunda capa:} un único elemento lógico de umbral fijo (OR, AND, MAJ).
  \end{itemize}
  \begin{center}
    $\mathbf{X} \to$ [Adaline$_1$, $\ldots$, Adaline$_n$] $\to$ Lógica Fija $\to \pm 1$
  \end{center}
\end{frame}

\begin{frame}{MRI — Ejemplo 1 (Sin error)}
  \begin{center}
    \begin{tabular}{c|c|c|c}
      Adaline & $s = \mathbf{W}^T\mathbf{X}$ & $\sgn(s)$ & \\
      \hline
      $W_1$ & $0.2$ & $+1$ & \\
      $W_2$ & $0.8$ & $+1$ & \\
      $W_3$ & $-0.1$ & $-1$ & $\to$ MAJ $\to \mathbf{-1}$ \\
      $W_4$ & $-0.3$ & $-1$ & \\
      $W_5$ & $-0.4$ & $-1$ & \\
    \end{tabular}
  \end{center}
  \begin{itemize}
    \item Salida deseada $d = -1$. \textbf{¡Correcto!} No se requiere adaptación.
  \end{itemize}
\end{frame}

\begin{frame}{MRI — Ejemplo 2 (Error detectado)}
  \begin{center}
    \textbf{Mismo $\mathbf{X}$, pero $d = +1$}
    \begin{tabular}{c|c|c|c}
      Adaline & $s = \mathbf{W}^T\mathbf{X}$ & $\sgn(s)$ & \\
      \hline
      $W_1$ & $0.2$ & $+1$ & \\
      $W_2$ & $0.8$ & $+1$ & \\
      $W_3$ & $-0.1$ & $-1$ & $\to$ MAJ $\to \mathbf{-1}$ \\
      $W_4$ & $-0.3$ & $-1$ & \\
      $W_5$ & $-0.4$ & $-1$ & \\
    \end{tabular}
  \end{center}
  \begin{itemize}
    \item Salida deseada $d = +1$, pero la salida es $-1$. \textbf{¡Error!}
    \item Se elige el Adaline con menor $|s_k|$: $W_3$ ($|s| = 0.1$, el más cercano a cero).
  \end{itemize}
\end{frame}

\begin{frame}{MRI — Ejemplo 3 (Adaptación de $W_3$)}
  \begin{center}
    \textbf{Se adapta $W_3$ para invertir su salida:}
    \begin{tabular}{c|c|c|c}
      Adaline & $s = \mathbf{W}^T\mathbf{X}$ & $\sgn(s)$ & \\
      \hline
      $W_1$ & $0.2$ & $+1$ & \\
      $W_2$ & $0.8$ & $+1$ & \\
      $\mathbf{W_3'}$ & $\mathbf{0.4}$ & $\mathbf{+1}$ & $\to$ MAJ $\to \mathbf{+1}$ \\
      $W_4$ & $-0.3$ & $-1$ & \\
      $W_5$ & $-0.4$ & $-1$ & \\
    \end{tabular}
  \end{center}
  \begin{itemize}
    \item ¡Mayoría ahora es $+1$! Salida deseada $d = +1$. \textbf{¡Correcto!}
  \end{itemize}
\end{frame}

\begin{frame}{MRI — Intento alternativo (perturbando $W_2$)}
  \begin{center}
    \textbf{Perturbando $W_2$:}
    \begin{tabular}{c|c|c|c}
      Adaline & $s = \mathbf{W}^T\mathbf{X}$ & $\sgn(s)$ & \\
      \hline
      $W_1$ & $0.2$ & $+1$ & \\
      $\mathbf{W_2'}$ & $\mathbf{-0.05}$ & $\mathbf{-1}$ & \\
      $W_3$ & $-0.1$ & $-1$ & $\to$ MAJ $\to \mathbf{-1}$ \\
      $W_4$ & $-0.3$ & $-1$ & \\
      $W_5$ & $-0.4$ & $-1$ & \\
    \end{tabular}
  \end{center}
  \begin{itemize}
    \item Salida deseada $d = +1$, pero salida sigue siendo $-1$. ¡No funciona!
    \item Se necesita adaptar más de un Adaline.
  \end{itemize}
\end{frame}

\begin{frame}{MRI — Adaptación de $W_2'$ y $W_3'$}
  \begin{center}
    \begin{tabular}{c|c|c|c}
      Adaline & $s = \mathbf{W}^T\mathbf{X}$ & $\sgn(s)$ & \\
      \hline
      $W_1$ & $0.2$ & $+1$ & \\
      $\mathbf{W_2'}$ & $\mathbf{0.1}$ & $\mathbf{+1}$ & \\
      $\mathbf{W_3'}$ & $\mathbf{0.4}$ & $\mathbf{+1}$ & $\to$ MAJ $\to \mathbf{+1}$ \\
      $W_4$ & $-0.3$ & $-1$ & \\
      $W_5$ & $-0.4$ & $-1$ & \\
    \end{tabular}
  \end{center}
  \begin{itemize}
    \item ¡Ahora sí! Mayoría $+1$, $d = +1$. \textbf{¡Correcto!}
  \end{itemize}
\end{frame}

\begin{frame}{MRI — Algoritmo}
  \begin{itemize}
    \item Los pesos se inicializan con valores aleatorios pequeños.
    \item Actualización del vector de pesos: puede cambiarse agresivamente usando corrección absoluta o adaptarse por el pequeño incremento determinado por $\alpha$-LMS.
    \item \textbf{Principio:} asignar responsabilidad al Adaline o Adalines que pueden asumirla más fácilmente.
    \item La secuencia de presentación de patrones debe ser aleatoria.
  \end{itemize}
\end{frame}

\begin{frame}{Madaline Rule II (MRII) — Arquitectura}
  \begin{itemize}
    \item Arquitectura típica de Madaline II de dos capas.
    \item Ambas capas son adaptativas.
  \end{itemize}
\end{frame}

\begin{frame}{Madaline Rule II (MRII) — Algoritmo}
  \begin{itemize}
    \item Pesos inicializados con valores aleatorios pequeños.
    \item Secuencia de presentación de patrones aleatoria.
  \end{itemize}

  \begin{block}{Adaptación de los Adalines de la primera capa}
    \begin{enumerate}
      \item Seleccionar la menor magnitud de salida lineal $|s_k|$.
      \item Realizar una ``adaptación de prueba'' agregando una perturbación $\Delta s$ de amplitud adecuada para invertir su salida binaria.
      \item Si el error de salida se reduce, remover $\Delta s$ y cambiar el peso del Adaline seleccionado usando $\alpha$-LMS.
      \item Perturbar y actualizar otros Adalines en la primera capa con salida $s_k$ ``suficientemente pequeña'' (se pueden invertir pares, tríos, etc.).
      \item Después de agotar posibilidades en la primera capa, pasar a la siguiente capa.
      \item Seleccionar aleatoriamente un nuevo patrón de entrenamiento y repetir.
    \end{enumerate}
  \end{block}
\end{frame}

% ═══════════════════════════════════════════════════════════════════════════
\section{Descenso más Empinado (Steepest Descent)}
% ═══════════════════════════════════════════════════════════════════════════

\begin{frame}{Descenso más Empinado — Concepto}
  \begin{itemize}
    \item \textbf{Descenso más empinado:} medir el gradiente de la función de error cuadrático medio y alterar el vector de pesos en la dirección negativa del gradiente medido.
  \end{itemize}

  \begin{equation*}
    \mathbf{W}_{k+1} = \mathbf{W}_k + \mu(-\nabla_k)
  \end{equation*}

  \begin{itemize}
    \item $\mu$: controla estabilidad y tasa de convergencia.
    \item $\nabla_k$: valor del gradiente en un punto de la superficie MSE correspondiente a $\mathbf{W} = \mathbf{W}_k$.
  \end{itemize}
\end{frame}

% ── Descenso más Empinado: Elemento Único ───────────────────────────────────

\begin{frame}{$\mu$-LMS — Elemento Único}
  \begin{itemize}
    \item Error lineal instantáneo: $\varepsilon_k = d_k - \mathbf{W}_k^T \mathbf{X}_k$.
    \item Gradiente instantáneo:
    $\hat{\nabla}_k = \dfrac{\partial \varepsilon_k^2}{\partial \mathbf{W}_k} = -2\varepsilon_k \mathbf{X}_k$
    \item Actualización de pesos:
    $\mathbf{W}_{k+1} = \mathbf{W}_k + \mu(-\hat{\nabla}_k) = \mathbf{W}_k + 2\mu\varepsilon_k \mathbf{X}_k$
  \end{itemize}
\end{frame}

\begin{frame}{$\mu$-LMS — Derivación del Gradiente}
  \begin{align*}
    \varepsilon_k^2 &= (d_k - \mathbf{X}_k^T \mathbf{W}_k)^2 = d_k^2 - 2d_k \mathbf{X}_k^T \mathbf{W}_k + \mathbf{W}_k^T (\mathbf{X}_k \mathbf{X}_k^T) \mathbf{W}_k \\[4pt]
    \nabla_k &= \frac{\partial \varepsilon_k^2}{\partial \mathbf{W}_k} = -2d_k \mathbf{X}_k + 2(\mathbf{X}_k \mathbf{X}_k^T) \mathbf{W}_k \\[4pt]
    &= -2\mathbf{X}_k (d_k - \mathbf{X}_k^T \mathbf{W}_k) = -2\varepsilon_k \mathbf{X}_k
  \end{align*}

  \begin{itemize}
    \item $\mu$ controla estabilidad y tasa de convergencia.
    \item $d_k$ y $\mathbf{X}_k$ son independientes de $\mathbf{W}_k$.
  \end{itemize}
\end{frame}

\begin{frame}{$\mu$-LMS — Comportamiento Estadístico}
  \begin{itemize}
    \item Con $\mu$ pequeño y $K$ pequeño (pocas presentaciones de entrenamiento):
    \begin{itemize}
      \item El cambio total de peso durante la presentación de $K$ patrones es proporcional a la derivada del error cuadrático medio.
    \end{itemize}
    \item $\varepsilon$ es la función de error cuadrático medio (MSE).
  \end{itemize}
\end{frame}

\begin{frame}{Reglas No Lineales — Adaline con Sigmoide}
  \begin{itemize}
    \item Se minimiza el cuadrado medio del error sigmoide $\tilde{\varepsilon}_k$:
    \begin{equation*}
      \tilde{\varepsilon}_k = d_k - y_k = d_k - \sgm(s_k)
    \end{equation*}
  \end{itemize}

  \begin{center}
    $s_k \to$ Sigmoide $\to y_k \quad\to\quad \tilde{\varepsilon}_k = d_k - y_k$
  \end{center}
\end{frame}

% ── Backpropagation y MRIII: Elemento Único ─────────────────────────────────

\begin{frame}{Backpropagation: Adaline con Sigmoide}
  \begin{itemize}
    \item Estimación del gradiente instantáneo:
  \end{itemize}
  \begin{align*}
    \tilde{\varepsilon}_k &= d_k - \sgm(s_k) \\
    \frac{\partial \tilde{\varepsilon}_k^2}{\partial \mathbf{W}_k}
    &= \frac{\partial \tilde{\varepsilon}_k^2}{\partial s_k} \cdot \frac{\partial s_k}{\partial \mathbf{W}_k} \\
    &= -2\tilde{\varepsilon}_k \sgm'(s_k) \mathbf{X}_k
  \end{align*}
  \begin{itemize}
    \item Actualización:
    $\mathbf{W}_{k+1} = \mathbf{W}_k + 2\mu\tilde{\varepsilon}_k \sgm'(s_k) \mathbf{X}_k$
  \end{itemize}
\end{frame}

\begin{frame}{MRIII: Adaline con Sigmoide}
  \begin{block}{Motivación}
    La función sigmoide y su derivada pueden no realizarse con precisión cuando se implementan con hardware analógico.
  \end{block}

  \begin{itemize}
    \item Se usa una perturbación en lugar de la derivada de la sigmoide.
    \item Se agrega una pequeña señal de perturbación $\Delta s$ a la suma $s_k$.
    \item Se observa el efecto de esta perturbación sobre la salida $y_k$ y el error $\tilde{\varepsilon}_k$.
  \end{itemize}
\end{frame}

\begin{frame}{MRIII — Estimación del Gradiente}
  \begin{align*}
    \frac{\partial \tilde{\varepsilon}_k^2}{\partial s_k} &\approx \frac{\tilde{\varepsilon}_k^2(s_k + \Delta s) - \tilde{\varepsilon}_k^2(s_k)}{\Delta s} = \frac{\Delta\tilde{\varepsilon}_k^2}{\Delta s} \\[4pt]
    \hat{\nabla}_k &= \frac{\partial \tilde{\varepsilon}_k^2}{\partial \mathbf{W}_k} \approx \frac{\Delta\tilde{\varepsilon}_k^2}{\Delta s} \mathbf{X}_k
  \end{align*}

  \begin{itemize}
    \item Alternativamente:
    \begin{equation*}
      \frac{\Delta\tilde{\varepsilon}_k^2}{\Delta s} \approx 2\tilde{\varepsilon}_k \frac{\Delta\tilde{\varepsilon}_k}{\Delta s}
    \end{equation*}
    \item Actualización de pesos:
    \begin{equation*}
      \mathbf{W}_{k+1} = \mathbf{W}_k - \mu \frac{\Delta\tilde{\varepsilon}_k^2}{\Delta s} \mathbf{X}_k
      \quad\text{o}\quad
      \mathbf{W}_{k+1} = \mathbf{W}_k - 2\mu\tilde{\varepsilon}_k \frac{\Delta\tilde{\varepsilon}_k}{\Delta s} \mathbf{X}_k
    \end{equation*}
  \end{itemize}
\end{frame}

\begin{frame}{Backpropagation vs MRIII (Elemento Único)}
  \begin{itemize}
    \item Matemáticamente equivalentes si la perturbación $\Delta s$ es pequeña.
    \item Agregar la perturbación $\Delta s$ causa un cambio en $\tilde{\varepsilon}_k$ igual a:
    \begin{equation*}
      \Delta\tilde{\varepsilon}_k = \tilde{\varepsilon}_k(s_k + \Delta s) - \tilde{\varepsilon}_k(s_k) \approx \frac{\partial \tilde{\varepsilon}_k}{\partial s_k} \Delta s
    \end{equation*}
    \item Cuando $\Delta s \to 0$:
    \begin{equation*}
      \frac{\Delta\tilde{\varepsilon}_k}{\Delta s} \to \frac{\partial \tilde{\varepsilon}_k}{\partial s_k} = -\sgm'(s_k)
    \end{equation*}
    \item MRIII $\equiv$ Backpropagation.
  \end{itemize}
\end{frame}

% ── Backpropagation y MRIII: Redes Multi-Elemento ───────────────────────────

\begin{frame}{Backpropagation para Redes}
  \begin{block}{¿Cómo funciona backpropagation?}
    Para un vector de patrón de entrada $\mathbf{X}$:
    \begin{enumerate}
      \item Barrer hacia adelante a través del sistema para obtener un vector de respuesta de salida $\mathbf{Y}$.
      \item Calcular los errores en cada salida.
      \item Barrer los efectos de los errores hacia atrás a través de la red para asociar una ``derivada del error cuadrático'' $\delta$ con cada Adaline.
      \item Calcular un gradiente a partir de cada $\delta$.
      \item Actualizar los pesos de cada Adaline basado en el gradiente correspondiente.
    \end{enumerate}
  \end{block}
\end{frame}

\begin{frame}{Backpropagation: Propagación hacia Adelante}
  \begin{align*}
    a^{(1)} &= \mathbf{X} \\
    a^{(2)} &= \sgm((W^{(1)})^T a^{(1)}) \\
    a^{(3)} &= \sgm((W^{(2)})^T a^{(2)}) \\
    a^{(4)} &= \sgm((W^{(3)})^T a^{(3)})
  \end{align*}
  \begin{center}
    \textit{Figuras tomadas del curso online de Andrew Ng}
  \end{center}
\end{frame}

\begin{frame}{Backpropagation: Retro-propagación del Error}
  \begin{itemize}
    \item Para cada unidad de salida en la capa 4:
    \begin{equation*}
      \delta_j^{(4)} = a_j^{(4)} - y_j
    \end{equation*}
    \item Retro-propagación:
    \begin{align*}
      \delta^{(3)} &= W^{(3)} \delta^{(4)} \;.\!*\; \sgm'(W^{(2)} a^{(2)}) \\
      \delta^{(2)} &= W^{(2)} \delta^{(3)} \;.\!*\; \sgm'(W^{(1)} a^{(1)})
    \end{align*}
  \end{itemize}

  \begin{block}{Intuición}
    $\delta_j^{(l)}$ = ``error'' del nodo $j$ en la capa $l$.
  \end{block}
  \begin{center}
    \textit{Figuras tomadas del curso online de Andrew Ng}
  \end{center}
\end{frame}

\begin{frame}{Backpropagation: Gradiente Instantáneo para un Adaline}
  \begin{itemize}
    \item Gradiente instantáneo para un elemento Adaline:
    \begin{equation*}
      \hat{\nabla}_k = \frac{\partial \varepsilon^2}{\partial \mathbf{W}_k}
    \end{equation*}
    \item Definiendo $\delta_k = \dfrac{\partial \varepsilon^2}{\partial s_k}$:
    \begin{equation*}
      \hat{\nabla}_k = \delta_k \mathbf{X}_k
    \end{equation*}
    \item Y la actualización:
    \begin{equation*}
      \mathbf{W}_{k+1} = \mathbf{W}_k - \mu \delta_k \mathbf{X}_k
    \end{equation*}
  \end{itemize}
\end{frame}

\begin{frame}{MRIII para Redes}
  \begin{itemize}
    \item Error cuadrático total de salida:
    \begin{equation*}
      \varepsilon^2 = (d_1 - y_1)^2 + (d_2 - y_2)^2 + \cdots
    \end{equation*}
    \item Se proporciona una perturbación agregando $\Delta s$ a un Adaline seleccionado; la perturbación se propaga a través de la red, causando un cambio en la suma de los cuadrados de los errores.
    \item Gradiente instantáneo:
    \begin{equation*}
      \hat{\nabla}_k = \frac{\partial \varepsilon_k^2}{\partial \mathbf{W}_k}
      \approx \frac{\Delta\varepsilon_k^2}{\Delta s} \mathbf{X}_k
    \end{equation*}
    \item Actualización de pesos:
    \begin{equation*}
      \mathbf{W}_{k+1} = \mathbf{W}_k - \mu \frac{\Delta\varepsilon_k^2}{\Delta s} \mathbf{X}_k
    \end{equation*}
  \end{itemize}
\end{frame}

% ═══════════════════════════════════════════════════════════════════════════
\section{Referencias}
% ═══════════════════════════════════════════════════════════════════════════

\begin{frame}{Referencias}
  \footnotesize
  \begin{itemize}
    \item B. Widrow and M.E. Hoff, Jr. \textit{Adaptive switching circuits}. Technical Report 1553-1, Stanford Electron. Labs., Stanford, CA, June 30, 1960.
    \item K. Senne. \textit{Adaptive Linear Discrete-time estimation}. PhD thesis, Technical Report SEL-68-090, Stanford Electron. Labs., Stanford, CA, June 1968.
    \item S. Grossberg. \textit{Nonlinear neural networks: principles, mechanisms, and architectures}. Neural Networks 1.1 (1988): 17--61.
    \item S. Cross, R. Harrison, and R. Kennedy. \textit{Introduction to neural networks}. The Lancet 346.8982 (1995): 1075--1079.
    \item Andrew Ng's Machine Learning course on Coursera.
    \item LMS tutorials: \url{http://homepages.cae.wisc.edu/~ece539/resources/tutorials/LMS.pdf}
  \end{itemize}
\end{frame}

\end{document}
